% =============================================================================
%  Capítulo 1 -- Introducción
% =============================================================================
\chapter{Introducción}
\label{cap:introduccion}

\section{Contexto del problema}
\label{sec:contexto}

Los sistemas de clasificación automática de texto basados en técnicas de
procesamiento de lenguaje natural (\PLN) se utilizan cada vez con mayor
frecuencia en procesos como el filtrado automatizado de currículums, la
moderación de contenido o la categorización de información en plataformas
digitales. Aunque estos sistemas suelen presentarse como herramientas neutrales
orientadas a la eficiencia, diversas investigaciones han mostrado que pueden
reproducir o amplificar desigualdades sociales preexistentes cuando aprenden
de datos históricos o de patrones lingüísticos asociados a contextos sociales
específicos. En particular, los modelos de \PLN pueden incorporar sesgos
provenientes de los datos, del proceso de anotación o del diseño del sistema,
lo que puede generar resultados sistemáticamente desfavorables para ciertos
grupos sociales \cite{HovySpruit2016,HovyPrabhumoye2021}.

Desde la sociolingüística se ha documentado que el uso del lenguaje varía de
acuerdo con factores como el nivel educativo, la clase social o el contexto
cultural, y que estas variaciones pueden ser objeto de estigmatización o
discriminación. La \emph{discriminación lingüística} ocurre cuando determinados
registros, dialectos o estilos de habla se asocian con juicios negativos sobre
la capacidad, la inteligencia o la credibilidad de una persona, reproduciendo
jerarquías sociales preexistentes \cite{Drozdzowicz2024,Weissler2022}. En
entornos digitales mediados por algoritmos, estas asociaciones podrían
trasladarse a sistemas automatizados que procesan lenguaje natural, generando
respuestas diferenciadas ante variaciones sociolectales que no afectan el
contenido sustantivo de la información.

\section{Planteamiento del problema}
\label{sec:planteamiento}

A pesar del creciente uso de sistemas de \PLN en procesos con consecuencias
sociales relevantes, la mayoría de las evaluaciones de equidad disponibles en
la literatura se han centrado en atributos demográficos categóricos
---típicamente género, raza o edad--- y en lenguas distintas del español. La
dimensión \emph{sociolectal}, entendida aquí como variación lingüística
asociada a factores socioeconómicos, ha recibido menos atención sistemática,
particularmente en el español hablado en México. Como consecuencia, no se
dispone de instrumentos estandarizados para evaluar si un clasificador de
texto responde de manera consistente cuando se le presenta el mismo contenido
informativo expresado en registros distintos.

Este vacío instrumental tiene implicaciones prácticas. En procesos como el
filtrado automatizado de currículums, un comportamiento sistemáticamente
distinto del modelo ante variaciones de registro puede traducirse en
desventajas concretas para personas cuyo estilo de escritura difiere del
registro dominante en los datos de entrenamiento, sin que ello refleje
diferencias reales en el contenido evaluado.

\section{Pregunta de investigación}
\label{sec:pregunta}

\begin{citablock}
¿En qué medida los sistemas de clasificación automatizada de texto presentan
sesgos asociados al registro lingüístico o sociolectal, y cómo puede diseñarse
un indicador que permita evaluar su neutralidad frente a estas variaciones?
\end{citablock}

\section{Objetivos}
\label{sec:objetivos}

\subsection*{Objetivo general}

Analizar la presencia de sesgos asociados al registro lingüístico o sociolectal
en sistemas de clasificación automatizada de texto, mediante la construcción de
textos controlados que varían únicamente en su estilo lingüístico, con el
propósito de evaluar su comportamiento frente a estas variaciones y proponer un
indicador denominado \emph{Índice de Neutralidad Sociolectal} (\INS) para medir
el grado de equidad lingüística en dichos sistemas.

\subsection*{Objetivos específicos}

\begin{enumerate}[label=\textbf{O\arabic*.},leftmargin=2.2em]
  \item Documentar, mediante revisión de literatura académica y reportes
    institucionales, cómo el clasismo estructural y la discriminación
    lingüística se manifiestan en el contexto mexicano y de qué manera podrían
    reflejarse en sistemas automatizados de procesamiento de lenguaje.
  \item Diseñar un conjunto de textos o currículums simulados en los que el
    contenido informativo permanezca constante, variando únicamente el
    registro lingüístico, con el fin de construir un conjunto de pruebas
    controladas para evaluar sistemas de clasificación automática de texto.
  \item Evaluar el comportamiento de uno o más modelos de clasificación
    textual frente a las variaciones lingüísticas definidas, comparando sus
    resultados para identificar posibles diferencias sistemáticas asociadas
    al estilo de escritura.
  \item Diseñar y calcular un indicador denominado \emph{Índice de Neutralidad
    Sociolectal} (\INS) que permita medir el grado de consistencia o variación
    en las respuestas del sistema ante cambios en el registro lingüístico.
\end{enumerate}

\section{Alcances y limitaciones}
\label{sec:alcances}

\paragraph{Alcances.} El estudio se restringe a sistemas de clasificación
automática de texto sobre lengua española, con énfasis en el español hablado
en México. La evaluación experimental se circunscribe a uno o más modelos de
clasificación seleccionados explícitamente en el Capítulo~\ref{cap:metodologia}
y al \emph{benchmark} controlado descrito en el Capítulo~\ref{cap:resultados}.

\paragraph{Limitaciones.} Las limitaciones principales del trabajo son tres.
Primero, el alcance muestral de la encuesta de percepciones no permite
inferencias representativas a nivel poblacional; sus resultados se interpretan
como insumo cualitativo para el diseño del \emph{benchmark}. Segundo, el
\emph{benchmark} controlado opera sobre textos simulados y no captura la
totalidad de la variación sociolectal observada en datos reales. Tercero, el
\INS se propone como indicador \emph{preliminar} y su validación empírica
exhaustiva se identifica como trabajo futuro.

\section{Contribuciones de la tesis}
\label{sec:contribuciones}

Este trabajo aporta los siguientes elementos.

\begin{itemize}
  \item Una articulación entre la literatura sobre discriminación lingüística
    y la de evaluación de sesgos en \PLN, situada en el contexto mexicano.
  \item Un \emph{benchmark} controlado de textos en distintos registros
    sociolectales, con contenido informativo equivalente.
  \item Un protocolo experimental reproducible para evaluar la consistencia
    de un clasificador frente a variaciones de registro.
  \item La definición operativa y el cálculo de un indicador, el \INS, sobre
    los resultados experimentales.
\end{itemize}

\section{Estructura del documento}
\label{sec:estructura}

El resto de la tesis se organiza como sigue. El Capítulo~\ref{cap:marco}
presenta el marco teórico, articulando los conceptos de sesgo algorítmico,
clasificación automática de texto, discriminación lingüística y clasismo
estructural en el contexto mexicano. El Capítulo~\ref{cap:metodologia}
detalla la metodología sociotécnica empleada. El Capítulo~\ref{cap:resultados}
reporta los resultados de la encuesta de percepciones y de los experimentos
computacionales. El Capítulo~\ref{cap:ins} define formalmente el \INS,
describe su procedimiento de cálculo y discute su interpretación. El
Capítulo~\ref{cap:discusion} discute las implicaciones sociotécnicas de los
resultados y las limitaciones metodológicas del estudio. El
Capítulo~\ref{cap:conclusiones} presenta las conclusiones generales y
delinea líneas de investigación futura.
